% Vietnamese translation for OI Public Library
% Source-File: IOI中国国家候选队论文/2007/day2/10.刘家骅《浅谈随机化在信息学竞赛中的应用》.ppt
\documentclass[11pt]{article}
\usepackage{oipl-vn}

\newcommand{\Oh}{\mathcal{O}}

\title{Bản ghi chú trình chiếu: Ứng dụng ngẫu nhiên hóa trong thi tin học}
\author{Liu Jiahua}
\date{2007}

\begin{document}
\maketitle

\origtitle{Bàn sơ về ứng dụng ngẫu nhiên hóa trong thi tin học}
\sourcepath{IOI China candidate team papers / 2007 / day 2 / Liu Jiahua.ppt}

\begin{abstract}
Các slide giới thiệu ngẫu nhiên hóa như một cách đổi hướng tìm kiếm khi lời
giải xác định khó nghĩ, khó cài đặt, hoặc không cần bảo đảm tối ưu tuyệt đối.
Qua ba nhóm ví dụ, bài nhấn mạnh rằng ngẫu nhiên hóa không phải là ``làm bừa'',
mà cần kết hợp với phân tích không gian trạng thái, tiền xử lý và cải thiện
cục bộ.
\end{abstract}

\paragraph{Từ khóa.}
Ngẫu nhiên hóa; tìm kiếm cục bộ; thu hẹp không gian tìm kiếm; đề mở; cây
khung có ràng buộc bậc.

\section{Vai trò của ngẫu nhiên hóa}

Trong thi tin học, thuật toán ngẫu nhiên hóa có thể xuất hiện dưới nhiều vai
trò:
\begin{itemize}
  \item sinh nghiệm ban đầu hoặc sinh hướng thử;
  \item phá cấu trúc dữ liệu đối kháng;
  \item thu hẹp vùng tìm kiếm trước khi vét cạn chính xác;
  \item cải thiện nghiệm bằng thao tác cục bộ lặp nhiều lần;
  \item tìm nghiệm đủ tốt cho đề mở hoặc bài không yêu cầu tối ưu tuyệt đối.
\end{itemize}

Điểm phải kiểm soát là xác suất thất bại và thời gian. Một lời giải ngẫu
nhiên tốt thường vẫn dựa trên cấu trúc của bài toán: có tiền xử lý để đánh
giá nhanh, có tiêu chí nhận hoặc từ chối bước ngẫu nhiên, và có điều kiện
dừng rõ ràng.

\section{Geometrical Dreams}

\subsection{Bài toán}

Cho đa giác \(A_1A_2\ldots A_N\). Trên mỗi cạnh \(A_iA_{i+1}\), dựng ra phía
ngoài một tam giác cân \(A_iM_iA_{i+1}\) sao cho
\(\angle A_iM_iA_{i+1}=\alpha_i\). Cho \(N\), các điểm \(M_i\) và các góc
\(\alpha_i\), cần khôi phục tọa độ các đỉnh \(A_i\). Tập các góc thỏa điều
kiện không có tổng của tập con không rỗng nào là bội của \(360^\circ\).

Bài này có lời giải chuẩn bằng hệ phương trình. Slide cố ý dùng nó như một
ví dụ nhỏ để giới thiệu một góc nhìn ngẫu nhiên hóa.

\subsection{Ý tưởng ngẫu nhiên hóa}

Nếu biết tọa độ của một đỉnh, các đỉnh còn lại có thể dựng tuần tự bằng các
phép hình học đơn giản. Vì vậy bài toán chuyển thành tìm tọa độ \(A_1\).
Đặt tạm \(A_1\) ở một vị trí nào đó, dựng lần lượt đến \(A_{N+1}\). Nếu
\(A_{N+1}=A_1\), ta đã tìm được đa giác đúng.

Định nghĩa hàm sai lệch là khoảng cách giữa điểm đầu và điểm cuối sau khi
dựng:
\[
  \Delta(A_1)=|A_{N+1}-A_1|.
\]
Ta chọn ngẫu nhiên một điểm gần vị trí hiện tại của \(A_1\). Nếu vị trí mới
làm \(\Delta\) nhỏ hơn thì nhận nó; nếu không thì bỏ. Lặp lại quá trình này
và giảm phạm vi thử khi cần, cho đến khi sai số đủ nhỏ.

Phương pháp này không phải cách hay nhất cho riêng bài này, nhưng nó cho thấy
một mẫu quan trọng: dùng ngẫu nhiên hóa để tìm cực tiểu của một hàm sai lệch,
khi chỉ cần đánh giá nhanh chất lượng của một trạng thái.

\section{Two Sawmills}

\subsection{Bài toán}

Trên đường từ đỉnh núi xuống chân núi có \(n\) cây. Gỗ chỉ được vận chuyển
theo chiều xuống dốc. Ở cuối đường đã có một xưởng cưa, và có thể xây thêm
hai xưởng tại vị trí của hai cây. Cần chọn vị trí sao cho tổng chi phí vận
chuyển nhỏ nhất.

Thuật toán chuẩn chuyển dữ liệu thành một bài hình học và dùng ngăn xếp để
tính trong \(\Oh(N)\), nhưng không dễ nghĩ. Slide dùng bài này để minh họa
ngẫu nhiên hóa có cấu trúc.

\subsection{Chọn ngẫu nhiên trực tiếp}

Cách đầu tiên là chọn ngẫu nhiên hai vị trí xưởng cưa, tính chi phí, rồi lặp
lại nhiều lần và lấy nghiệm tốt nhất. Với tiền xử lý, chi phí của mỗi cặp có
thể tính trong \(\Oh(1)\). Tuy vậy, nếu \(n\) tới \(20000\), không gian trạng
thái có khoảng \(4\cdot 10^8\) cặp; chọn ngẫu nhiên trực tiếp vài triệu lần
vẫn không bảo đảm gặp nghiệm tối ưu.

\subsection{Thu hẹp vùng tìm kiếm}

Cách tốt hơn là xem mỗi cặp vị trí \((X,Y)\) là một điểm trong ma trận \(P\),
trong đó \(P[X,Y]\) là tổng chi phí. Ban đầu vùng tìm kiếm là toàn bộ ma trận
\(N\times N\). Ở mỗi vòng:
\begin{enumerate}
  \item chọn ngẫu nhiên một số điểm trong vùng hiện tại;
  \item tính giá trị \(P\) tại các điểm đó và lấy điểm tốt nhất làm tâm;
  \item cắt ra một vùng con quanh tâm, chẳng hạn cạnh bằng \(3/4\) cạnh hiện
  tại;
  \item nếu vùng con vượt biên, dịch nó vào trong vùng cũ;
  \item lặp lại cho đến khi vùng đủ nhỏ, rồi vét cạn chính xác trong vùng đó.
\end{enumerate}

Ngẫu nhiên hóa ở đây không trực tiếp ``đoán đáp án'', mà dùng để nhanh chóng
tìm khu vực đáng xét. Phần cuối vẫn là vét cạn xác định trên vùng nhỏ, nên độ
ổn định cao hơn nhiều.

\section{Buổi tụ họp của Tiểu H}

\subsection{Bài toán}

Đây là bài nộp đáp án NOI 2005. Có \(N\) người bạn và các cặp có thể liên lạc,
mỗi cạnh có mức độ vui vẻ. Cần chọn \(N-1\) cạnh tạo thành một mạng liên thông
với tổng vui vẻ càng lớn càng tốt, đồng thời mỗi đỉnh \(i\) có giới hạn bậc
\(k_i\).

Khi mọi \(k_i=N-1\), bài toán là cây khung lớn nhất. Khi chỉ một đỉnh có giới
hạn bậc, có thể dùng thuật toán cây khung có ràng buộc bậc đặc biệt. Khi mọi
đỉnh đều bị giới hạn, bài toán trở thành khó kiểu NP; hơn nữa đây là đề mở,
không yêu cầu chứng minh tối ưu tuyệt đối. Đây là môi trường phù hợp cho
ngẫu nhiên hóa kết hợp cải thiện cục bộ.

\subsection{Phương pháp dựa trên Kruskal}

Sắp cạnh theo trọng số giảm dần, nhưng ngẫu nhiên xáo trộn một phần thứ tự
cạnh. Chạy Kruskal với điều kiện chỉ thêm cạnh nếu không tạo chu trình và
không làm vượt giới hạn bậc. Vì thứ tự cạnh thay đổi, mỗi lần chạy có thể cho
một cây khác nhau.

Sau khi có một cây ứng viên, chọn ngẫu nhiên một cạnh ngoài cây và thêm vào
cây để tạo một chu trình. Trên chu trình đó, thử xóa một cạnh sao cho cây vẫn
liên thông, vẫn thỏa giới hạn bậc, và tổng trọng số tăng. Lặp lại bước cải
thiện cục bộ cho đến khi không còn cải thiện rõ ràng, rồi chạy lại từ nhiều
thứ tự ngẫu nhiên khác nhau. Kết quả cuối là nghiệm tốt nhất trong các lần
thử.

\subsection{Phương pháp dựa trên Prim}

Trước hết tính một cây khung lớn nhất khi bỏ qua giới hạn bậc. Cây này có
tổng trọng số cao nhưng có thể vi phạm nhiều giới hạn. Sau đó lặp:
\begin{enumerate}
  \item chọn ngẫu nhiên một đỉnh đang vượt giới hạn bậc;
  \item tìm một cạnh ngoài cây có thể thay thế một cạnh trong cây nối với
  đỉnh đó;
  \item phép thay thế phải giảm bậc đỉnh vi phạm và giữ cây liên thông;
  \item ưu tiên thay thế làm giảm tổng trọng số ít nhất hoặc cải thiện toàn
  cục tốt nhất.
\end{enumerate}
Nếu sửa được đến khi mọi đỉnh thỏa giới hạn, lưu nghiệm; nếu mắc kẹt, khởi
động lại từ cây ban đầu hoặc một biến thể ngẫu nhiên khác.

\subsection{So sánh}

Hai cách dùng ngẫu nhiên hóa khác nhau: Kruskal ngẫu nhiên hóa quá trình xây
cây ngay từ đầu, còn Prim bắt đầu bằng nghiệm rất tốt theo trọng số rồi sửa
các vi phạm. Bảng kết quả trong bài gốc cho thấy hai cách có điểm mạnh khác
nhau trên từng bộ dữ liệu. Kết luận thực dụng là không nên chỉ viết một phép
random đơn lẻ; với đề mở, nên kết hợp nhiều chiến lược và giữ nghiệm tốt nhất.

\section{Kinh nghiệm sử dụng}

Các slide rút ra ba lưu ý chính:
\begin{itemize}
  \item Dùng ngẫu nhiên hóa linh hoạt theo cấu trúc bài, không thay thế hoàn
  toàn phân tích thuật toán.
  \item Kết hợp ngẫu nhiên hóa với thuật toán xác định thích hợp: Kruskal,
  Prim, tiền xử lý truy vấn \(\Oh(1)\), hoặc vét cạn vùng nhỏ.
  \item Tận dụng giới hạn thời gian một cách có kiểm soát: số lần thử, kích
  thước vùng con, điều kiện nhận bước mới và điều kiện dừng đều phải rõ ràng.
\end{itemize}

Ngẫu nhiên hóa có ưu thế rõ trong bài tối ưu khó, bài mở và bài có không gian
trạng thái lớn nhưng đánh giá trạng thái nhanh. Tuy nhiên, trong bài yêu cầu
đáp án chính xác, cần phân tích xác suất sai hoặc dùng ngẫu nhiên hóa chỉ như
một bước hỗ trợ cho phần xác định.

\end{document}
